\scanpage{61}
Le fait que cela donne ce qu'il faut du point de vue cohomologique se déduit par exemple de 1.6. en suivant la construction précédente du point de vue cohomologique. Bien entendu, $u$ dépend linéairement de $z$, et si $z$ est num. \unsure{équi.|équiv.} à $0$, il en est de $\unsure{\overline{\mathrm m}|=}$ de $C$, donc de $u$ en vertu de 1.1.

Variante de la démonstration. On cherche $B_X\unsure{\mathord{-}|\to} A_Y\simeq B_Y^*$, i.e. une \unsure{clas.|classe|cl.} de corr. div. sur $B_X\times B_Y$. On la construit en termes de $\Theta_X\in A^1(X\times B_X)$, $\Theta_Y\in A^1(Y\times B_Y)$, $Z\in A^{n+m-1}(X\times Y)$, \deleted{par le procédé en \unsure{y} considérant $\unsure{\pi}:X\times B_X\times\inserted{\Theta_X}\unsure{Y}\times B_Y$ et} en formant la classe de \unsure{correspondance|correspondances} $\Theta_Y\circ Z\circ{}^t\Theta_X$ sur $B_X\times B_Y$,
\[
\ldots\,-
\]

\deleted{\uline{Corollaire} \inserted{2.2} Tout homomorphisme $B_X\mathord{-}A_Y$ est défini par un $Z_1\in\operatorname{Cyc}_1(X\times Y)$.}
% archive sheet number: 60

